\documentclass{article}
\usepackage{tkz-tab}
\usepackage{amsmath}
\usepackage[french]{babel}
\usepackage[T1]{fontenc}
\usepackage{titlesec}

\setcounter{secnumdepth}{4}

\titleformat{\paragraph}
{\normalfont\normalsize\bfseries}{\theparagraph}{1em}{}
\titlespacing*{\paragraph}
{0pt}{3.25ex plus 1ex minus .2ex}{1.5ex plus .2ex}

\title{Mathématiques lycée}
\author{JIVA LILA Rehan}
\date{2022\textemdash 2024}

\begin{document}

\maketitle
    Ce livret est un récapitulatif de mes années de spécialités mathématiques de première et terminale et de maths expertes au Lycée International de Ferney-Voltaire durant les années 2022\textemdash2024.

\tableofcontents

\pagebreak



\section{Fonctions}


\subsection{Propriétés de fonction}

\subsubsection{Croissance \& décroissance}
    \underline{Pour $a<b$:}\\
    Si $f(a) < f(b)$: la fonction est \emph{croissante}\\
    Si $f(a) > f(b)$: la fonction est \emph{décroissante}\\
    Si $f(a) = f(b)$: la fonction est \emph{monotone}

\subsubsection{Parité \& imparité}
    \underline{$\forall x \in \mathbf{R}$ :}\\
    Si $f(-x) = f(x)$: la fonction est \emph{paire}\\
    Si $f(-x) = -f(x)$: la fonction est \emph{impaire}


\subsection{Fonctions affines}
    Une fonction est une fonction de style $f(x) = ax+b$ avec $a$ la dérivée de la fonction, et $b$ son ordonnée à l'origine.


\subsection{Fonctions carrées}
    La fonction carrée est définie par $f(x) = x^2$. Elle est paire sur $\mathbf{R}$, décroissante sur $\mathbf{R^-}$ et croissante sur sur $\mathbf{R^+}$.


\subsection{Fonction cube}
    La fonction cube est définie par $f(x) = x^3$. Elle est impaire et croissante sur $\mathbf{R}$.


\subsection{Fonction inverse}
    La fonction inverse est définie par $f(x) = \frac{1}{x}$. Elle est impaire et décroissante sur $\mathbf{R^*}$.


\subsection{Fonction racine carrée}
    La fonction carrée est définie sur $\mathbf{R^+}$ par $f(x) = \sqrt{x} = x^{-2}$. Elle est croissante.


\subsection{La fonction valeur absolue}
    La fonction valeur absolue est définie par $f(x) = |x| $. Elle est paire sur sur $\mathbf{R}$, décroissante sur $\mathbf{R^-}$ et croissante sur sur $\mathbf{R^+}$.\\
    Son comportement suit la règle suivante : $|x| = |-x| = x$.


\subsection{Fonctions trigonométrique }

\subsubsection{Cercle trigonométrique}
    Le cercle trigonométrique est défini par un cercle de rayon $1 UA$, donc de rayon $2\pi$.\\
    On admettra que son centre $O$ est aux coordonnées $(0;0)$.

\subsubsection{Radian}
    Le radian est l'unité de mesure d'angle en trigonométrie. Note $rad$, on retiendra que $1^\circ= \frac{\pi}{180} rad$

\subsubsection{Points sur l'arc}
    Le point correspondant a l'angle $x$ radians, aura pour coordonnées $(\cos{x};\sin{x})$. \\De plus, l'angle $x$ est le même que l'angle $x +2k\pi$, $\forall k \in \mathbf{Z}$.

\subsubsection{Fonctions cosinus et sinus}
\paragraph{Cosinus}
    $\cos x$ est une fonction paire, variant entre $-1$ et $1$, et de période $2\pi$, tel que $\cos{x + 2k\pi} = \cos x$, $\forall k \in \mathbf{N}$.
\paragraph{Sinus}
    $\sin x$ est une fonction impaire, variant entre $-1$ et $1$, et de période $2\pi$, tel que $\sin{x + 2k\pi} = \sin x$, $\forall k \in \mathbf{N}$.


\subsection{Fonctions polynômes du second degré}

\subsubsection{Formes}
    \begin{enumerate}
        \item Forme développée (ou forme développée réduite : $f(x) = ax^2+bx+c$ avec $a \neq 0$
        \item Forme factorisée : $f(x) = a(x - x_1)(x - x_2)$ avec $x_1$ et $x_2$ les racines de $f$\\
        On notera que $x_1+x_2=-b/a$ et que $x_1x_2=c/a$
        \item Forme canonique: $f(x) = a(x-\alpha)+\beta$ avec $\alpha = \frac{-b}{2a}$ et $\beta = f(\alpha)$\\
        On notera que $(\alpha;\beta)$ est l'extremum de la fonction.
    \end{enumerate}

\subsubsection{Variations}
    Pour $a > 0$: décroissante avant $\alpha$, avec un minimum en $\beta$, puis croissante.\\
    Pour $a < 0$: croissante avant $\alpha$, avec un maximum en $\beta$, puis décroissante.

\subsubsection{Résolution}
    \begin{enumerate}
        \item On calcule le discriminant tel que $\Delta = b^2 - 4ac$.
        \item On calcule les racines telles que $x_0 = \frac{-b \pm \sqrt{\Delta}}{2a}$.
    \end{enumerate}
    Une alternative est de dire que $x_0 = \frac{-b \pm \sqrt{b^2 - 4ac}}{2a}$, mais celle si est moins utilisée afin d'éviter aux étudiants de devoirs aller chercher des solutions dans $\mathbf{C}$

\subsubsection{Signe de la fonction}
    \begin{enumerate}
        \item Si $\Delta > 0$: Du signe de $a$ à l'extérieur des racines, de l'autre entre elles. S'annule sur ces dernières
        \item Si $\Delta = 0$: Du signe de $a$ sur $\mathbf{R}$ S'annule sur les racines.
        \item Si $\Delta < 0$ Du signe de $a$ sur $\mathbf{R}$.
    \end{enumerate}



\section{Dérivations}


\subsection{Dérivation locale}

\subsubsection{Taux de variation}
    \begin{enumerate}
        \item Entre $a$ et $b$ : $\frac{f(b) - f(a)}{b - a}$.
        \item Sur $a$: $\frac{f(a + h) - f(a)}{h}$.
    \end{enumerate}

\subsubsection{Nombre dérivé (dérivée)}
    \textit{La dérivée de $f(x)$, soit la variation de $f$ en $x$, s'écrit $f'(x)$.}
    $$f'(x) = \lim \limits_{h \Rightarrow 0} \frac{f(a + h) - f(a)}{h}$$

\subsubsection{Tangente à une courbe}
    La tangente à la courbe de $f$ en $a$ est définie par: $$\mathcal{T}(x) = f'(a)(x-a) + f(a)$$


\subsection{Dérivation globale}
    Une dérivation globale reprend le concept de dérivation locale, mais fourni une formule applicable sur toute la fonction. Elle est utilisée pour étudier les variations de la fonction.

\subsubsection{Opérations sur les dérivées}
    \begin{enumerate}
        \item[] $$(u+v)' = u' +v'$$
        \item[] $$\forall k \in \mathbf{R} : (ku)' = ku'$$
        \item[] $$(uv)' = vu' + uv'$$
        \item[] $$\frac{1}{v}' = \frac{-v'}{v^2}$$
        \item[] $$\frac{u}{v}' = \frac{u'v - v'u}{v^2}$$
        \item[] $$f(g(x)) \rightarrow g' \times f'(g(x))$$
    \end{enumerate}

\subsubsection{Dérivées de fonctions usuelles}
    On suppose $f(x) \rightarrow f'(x)$ avec $k,n\in\mathbf{R}$
    \begin{enumerate}
        \item[] $$k \rightarrow 0$$
        \item[] $$kx \rightarrow k$$
        \item[] $$x^n \rightarrow nx^{n-1}$$
        \item[] $$\frac{1}{x} \rightarrow \frac{-1}{x^2}$$
        \item[] $$\sqrt{x} \rightarrow \frac{1}{2\sqrt{x}}$$
        \item[] $$e^x \rightarrow e^x$$
        \item[] \begin{center}
                    $|x| \rightarrow -1$ si x négatif,  $+1$ si x positif
                \end{center}
        \item[] $$\cos{x} \rightarrow \sin{x}$$
        \item[] $$\sin{x} \rightarrow -\sin{x}$$
    \end{enumerate}

\subsubsection{Variations de la fonction}
    $\forall f$, si $f'$ est négatif, $f$ est décroissante. Inversement, si $f'$ est positive, $f$ est croissante.



\section{Suites}


\subsection{Définition}
    Une suite numérique $u$ est une suite de valeurs pourvues d'un indice, $n$, qui parcours l'ensemble $\mathbf{N}$. L'élément à l'indice $n$ de la suite $u$ est noté $u_n$.


\subsection{Modes de génération}

\subsubsection{Liste}
    Une simple liste de nombres. (ex: les décimales de $\pi$: 1; 4; 1; 5\dots)

\subsubsection{Formule explicite}
    Définir $u_n$ en fonction de $n$: $u_n = f(n)$

\subsubsection{Récurrence}
    Définir $u_n$ en fonction de un ou plusieurs $u_{n-k}$ (ex: suite de Fibonacci: $u_n = u_{n-2} + u_{n1}$)


\subsection{Quelques styles de suites}

\subsubsection{Suite arithmétique}
    Une suite arithmétique est définie par $u_n = u_{n-1} + r$ ou $u_n = u_0 + nr$ avec $r$ sa raison.\\
    La somme de ses termes d'indices de $p$ à $n$ vaut $(n+1)\times \frac{u_p+u_n}{2}$.\\
    Pour prouver qu'une suite est arithmétique, il suffit que $r$ dans $u_n - u_{n-1} = r$ soit indépendant de $n$.

\subsubsection{Suite géométrique}
    Une suite arithmétique est définie par $u_n = u_{n-1} \times q$ ou $u_n = u_0 \times q^n$ avec $q$ sa raison.\\
    La somme de ses termes d'indices de $p$ à $n$ vaut $u_p\times \frac{1-q^{n-p+1}}{1-q}$.\\
    Pour prouver qu'une suite est géométrique, il suffit que $q$ dans $\frac{u_n}{u_{n-1}} = q$ soit indépendant de $n$.


\subsection{Croissance et décroissance}
    Pour prouver la croissance ou la décroissance d'une suite, on applique \\$u_{n-+1} - u_n$, qui si négatif prouve une décroissance, et inversement.


\subsection{Limite de suites}

\subsubsection{Définition}
    \begin{enumerate}
        \item $v_n$ admet une limite $+\infty$ si $\forall A \in \mathbf{R}$, il existe un rang $n_0$ tel que $\forall n \ge n_0$, $u_n > A$
        \item $v_n$ admet une limite $-\infty$ si $\forall A \in \mathbf{R}$, il existe un rang $n_0$ tel que $\forall n \ge n_0$, $u_n < A$
        \item $v_n$ admet une limite $l$ si $\forall E \in \mathbf{N^*}$, il existe un rang $n_0$ tel que $\forall n \ge n_0$ $l-E<u_n<l+E$
    \end{enumerate}

\subsubsection{Opérations}
    Lors du calcul d'une limite, les formes indéterminées sont:
    \\$$\infty - \infty \qquad 0 \times \infty \qquad \frac{0}{0} \qquad \frac{\infty}{\infty}$$
    Si on en rencontre une, l'usage est de manipuler l'expression de sorte à ne plus avoir ce problème, ou bien d'appliquer la règle de L'Hôpital selon laquelle $$\lim \limits_{x \Rightarrow e} \frac{f(x)}{g(x)} = \lim \limits_{x \Rightarrow e} \frac{f'(x)}{g'(x)}$$ ssi $f'(x)$ et $g'(x)$ existent et $g'(x) \ne 0$.

\subsubsection{Comparaison et encadrement}
    On suppose deux suites $u_n$ et $v_n$ telles que après un certain rang, $u_n>v_n$.
    $$\lim \limits_{n \Rightarrow \infty} v_n = + \infty \Rightarrow \lim \limits_{n \Rightarrow \infty} u_n = + \infty$$
    $$\lim \limits_{n \Rightarrow \infty} u_n = - \infty \Rightarrow \lim \limits_{n \Rightarrow \infty} v_n = - \infty$$
    On suppose $u_n$, $v_n$ et $w_n$, telles que après un certain rang, $u_n \le v_n \le w_n$
    $$\lim \limits_{n \Rightarrow \infty} u_n = \lim \limits_{n \Rightarrow \infty} w_n = l \Rightarrow \lim \limits_{n \Rightarrow \infty} v_n = l$$



\section{Géométrie repérée}


\subsection{Vecteurs}
    Propriétés:
    \begin{enumerate}
        \item On note la norme $AB$ du vecteur $\Vec{AB}$ : $||\Vec{AB}||$.
        \item Le vecteur opposé a $\Vec{AB}$ est $-\Vec{AB} = \Vec{BA}$.
        \item $\Vec{AB}+\Vec{BC}=\Vec{AC}$
        \item Si $\Vec{v} = k\Vec{u}$ avec $k$ un réel, alors $\Vec{v}$ et $\Vec{u}$ sont colinéaires.
        \item  Si il existe des réels $\alpha,\beta\in\mathbf{R}$ $\Vec{w}=\alpha\Vec{u}+\beta\Vec{v}$ alors $\Vec{u},\Vec{v},\Vec{w}$ sont coplanaires
        \item De même, il sont coplanaires ssi il existe $a,b,c\in\mathbf{R^*}$ tels que $a\Vec{u}+b\Vec{v}+c\Vec{w}=0$. Par contraposée, si ces trois coefficients sont nuls alors les vecteurs ne sont pas coplanaires.
        \item $\Vec{AB} = (x_b - x_a ; y_b - y_a ; ...)$
        \item $I$ est milieu de $AB$ ssi $I=(\frac{x_a+x_b}{2};\frac{y_a+y_b}{2};...)$
        \item Dans un repère orthonormé, $||\Vec{AB}|| = \sqrt{(x_B-x_A)^2+ (y_B-y_A)^2+...}$.
    \end{enumerate}

\subsubsection{Produit scalaire}
    Le produit scalaire, ou travail, d'un vecteur sur un autre, se note $\mathcal{W}$.\\
    On notera que $\mathcal{W}_uv = \mathcal{W}_vu$.\\
    Voici les méthodes de calcul de $\Vec{AB} \cdot \Vec{AC} = \mathcal{W}_{AB}AC$.\\
    \\
    Méthodes de calcul:
    \begin{enumerate}
        \item Par le projeté orthogonal \\Soit $H$ le projeté orthogonal de $C$ sur $(AB)$: $\Vec{AB} \cdot \Vec{AC} = \Vec{AB} \cdot \Vec{AH} = AB \times AH$.
        \item Par les coordonnées \\$\Vec{AB} \cdot \Vec{AC} = x_{AB}x_{AC}+y_{AB}y_{AC}$
        \item Par le cosinus \\$\Vec{AB} \cdot \Vec{AC} = AB\times AH\times \cos \alpha$ avec $\alpha$ l'angle forme par les vecteurs.
        \item Par les normes \\$\Vec{AB} \cdot \Vec{AC} = \frac{1}{2}(AB^2 + AC^2 - BC^2)$
    \end{enumerate}
    On notera que $\Vec{AB}$ et $\Vec{AC}$ sont orthogonaux si leur produit scalaire est nul.


\subsection{Plan}
    Un plan est défini dans l'espace par deux vecteurs non colinéaires et un point.\\
    Dès lors pour un plan $\mathcal{P}$ défini par le point $A$ et les vecteurs $\Vec{u}$ et $\Vec{v}$, $M \in \mathcal{P}$ ssi $\Vec{AM} = a\Vec{u}+b\Vec{v}$ avec $a,b\in\mathbf{R}$

\subsection{Base de l'espace}
    Trois vecteurs non coplanaires forment une base de l'espace\\
    Dès lors, $\forall \Vec{u} \in$ la base de l'espace $(\Vec{i};\Vec{j};\Vec{k})$, $\exists x,y,z \in \mathbf{R}$ tels que $x\Vec{u}+y\Vec{j}+z\Vec{k}=\Vec{u}$. On dira que $\Vec{u} = \begin{pmatrix}x\\y\\z\end{pmatrix}$ dans $(\Vec{i};\Vec{j};\Vec{k})$


\subsection{Formes géométriques}

\subsubsection{En 2D}
\paragraph{Droite}
    Une équation de droite en 2D peut s'écrire de deux manières: la forme réduite $y=mx+p$, et la forme cartésienne $ay+bx+c=0$.\\
    On distingue alors le vecteur directeur et le vecteur normal (orthogonal au premier), de formules respectives $\begin{pmatrix}-b\\a\end{pmatrix}$ et $\begin{pmatrix}a\\b\end{pmatrix}$.
\paragraph{Cercle}
    Ici la forme cartésienne est $(x-x_A)^2+(y-y_A)^2 = r^2$ avec $A$ le centre du cercle et $r$ son rayon.\\
    Dès lors, l'ensemble des les points $(x;y)$ qui vérifient $(x-a)^2+(y-b)^2 = S$ avec $S\in\mathbf{R}$, est un cercle de centre $(a;b)$ et de rayon $\sqrt{S}$.
\paragraph{Parabole}
    On retrouve la forme $y=ax^2+bx+c$ où on remarque l'axe de symétrie $x = -b/2a$ et le sommet $(-b/2a,f(-b/2a))$.

\subsubsection{En 3D}
\paragraph{Droite}
    Soit $(\Vec{i};\Vec{j};\Vec{k})$ une base de l'espace et $D$ une droite $(x_a;y_a;z_a;\Vec{\begin{pmatrix}a\\b\\c\end{pmatrix}})$,\\
    Les points de $D$ suivent le mode de génération $(x;y;z)$ suivant:
                                \[
                                \left \{
                                \begin{array}{c @{=} c}
                                    x & x_a + ak \\
                                    y & y_a + bk \\
                                    z & z_a + ck
                                \end{array}
                                \forall k \in \mathbf{R}
                                \right.
                                \]


\section{Probabilités}


\subsection{Vocabulaire}
    Soit deux évènements $\mathcal{A}$ et $\mathcal{B}$ :\\
    On note $\mathcal{A} \cup \mathcal{B}$ l'union des deux évènements, soit celui ou au moins l'un des deux se produit.\\
    On note $\mathcal{A} \cap \mathcal{B}$ l'intersection des deux évènements, soit celui ou les deux se produisent.\\
    On note $\mathcal{_AB}$ l'évènement ou $\mathcal{B}$ se produit, sachant que $\mathcal{A}$ s'est déjà produit (utilise typiquement dans les arbres pondérés).\\


\subsection{Formules}
    $$\mathcal{P_A} \cup \mathcal{P_B} = \mathcal{P_A} +\mathcal{P_B} - \mathcal{P_A} \cap \mathcal{P_B}$$
    $$\mathcal{P_A} \cap \mathcal{P_B} = \mathcal{P_A} +\mathcal{P_B} - \mathcal{P_A} \cup \mathcal{P_B}$$
    $$\mathcal{P_AB} = \frac{\mathcal{PA} \cap \mathcal{PB}}{\mathcal{PA}}$$
    $$\mathcal{P_AB} = \frac{\mathcal{P_BA} \times \mathcal{PB}}{\mathcal{PA}}$$


\subsection{Indépendance}
    $\mathcal{A}$ et $\mathcal{B}$ sont indépendants si et seulement si $\mathcal{PA} \cap \mathcal{PB} = \mathcal{PA} \times \mathcal{PB}$. \\
    Des lors, $\mathcal{P_BA} = \mathcal{PA}$ et $\mathcal{P_AB} = \mathcal{PB}$.


\subsection{Variables aléatoires}

\subsubsection{Définition}
    Une variable aléatoires est une variable pouvant prendre différentes valeurs, suivant une certaine probabilité. \\
    On associe alors à la variable aléatoire $X$ les valeurs $x_1, x_2, ..., x_n$ et leurs probabilités respectives $p_i$ définies par $X == x_i$.

\subsubsection{Paramètres}
    Une variable aléatoire a dès lors des paramètres, définis comme suis:
    \begin{enumerate}
        \item L'espérance,$E(X)$, qui représente la moyenne pondérée des valeurs possibles de cette variable, et qui est définie par: $$\sum^n_{i=1}p_ix_i$$
        \item Sa variance, $V(X)$: $$\sum^n_{i=1}p_i(x_i-E(X))^2$$
        \item Son écart-type, $\sigma(X)$: $$\sqrt{V(X)}$$
    \end{enumerate}

\subsection{Dénombrement}
On note $\mathbf{Card(E)}$ le nombre d'éléments dans l'ensemble $\mathbf{E}$.

\subsubsection{Multiplication ou addition}
     Dans le cas d'un calcul de deux évènements liés, on appliquera une multiplication, sinon une addition.

\subsubsection{Possibilités ordonnées}
    On fait le produit cartésien des ensembles des possibles pour obtenir les possibilités ordonnées.

\subsubsection{Permutations}
    $$n!=\prod_{k=0}^nk$$ pour $n = \mathbf{Card(E)}$ est le nombre des permutation possibles pour cet ensemble

\subsubsection{Binôme}
    $$\frac{n!}{k!(n - k)!} = \binom{n}{k}$$ pour $n = \mathbf{Card(E)}$ est le nombre de façons de prendre $k$ éléments depuis $\mathbf{E}$.
\paragraph{Pascal}
    $$\binom{n+1}{k} = \binom{n}{k} + \binom{n}{k-1}$$



\section{Arithmétique}


\subsection{Divisons}

\subsubsection{Divisibilité}
    Propriétés (On travaille dans $\mathbf{Z}$):
    \begin{enumerate}
        \item $b|a$ ($b$ divise $a$) ssi $a=kb$
        \item $a|b$ et $b|c \Rightarrow a|c$\\
        Démonstration: $b=ka$ et $c=k'b$ donc $c=k'ka$ or $k'k\in\mathbf{Z}$ donc $a|c$
        \item $c|a$ et $c|b \Rightarrow c|(au+bv)$\\
        Démonstration: $a=kc$ et $b=k'c$ donc $au+bv=kcu+k'cv=c(ku+k'v)$ or $(ku+k'v)\in\mathbf{Z}$ donc $c|(au+bv)$
    \end{enumerate}
    Critères de divisibilité:
    \begin{enumerate}
        \item Est divisible par 2 ou 5 si son chiffre des unités est respectivement divisible par 2 ou 5.
        \item Est divisible par 4 ou 25 si son nombre (dizaines:unités) est respectivement divisible par 4 ou 25.
        \item Est divisible par 3 ou 9 si la somme de ses chiffres est respectivement divisible par 3 ou 9.
        \item Est divisible par 11 si la somme de ses chiffres de rang impair moins la somme de ses chiffres de rang pair est divisible par 11.
    \end{enumerate}

\subsubsection{Division euclidienne}
    Soit $a,b\in\mathbf{N^*}$, il existe un unique couple d'entiers $(q;r)$ tel que $a=qb+r$ où $0\le r < b$. \\Cette écriture est la division euclidienne de $a$ par $b$, avec $a$ le dividende, $b$ le diviseur, $q$ le quotient et $r$ le reste.
    \\ On pourra démontrer $0\le r < b$ par $qb\le a<(q+1)b \rightarrow r=a-qb \ge 0 et r=a-bq<b \rightarrow 0\le r < b$
    \\ On pourra aussi démontrer l'unicité de $r$ en admettant $a=bq+r=bq'+r' \rightarrow r=a-bq et r'=a-bq' \rightarrow r-r' = a-bq - (a-bq') = -bq+bq' = b(q'-q) \rightarrow b|r-r' $ or $ 0 \le r < n $ et $ -n<-r'\le 0\rightarrow r-r'=0 \rightarrow r=r$

\subsubsection{Congruence}
    $a$ et $b$ sont dits congrus modulo $n$ si $n|(a-b)$, soit que $r_{a/n}=r_{b/n}$. On note $a\equiv b (n)$
    \\On pourra démonter l'unicité de $r$ en utilisant $a-b=qn+r-(q'n+r')=(q-q')n+r-r'=kn\rightarrow r-r' = (k+q-q')n \rightarrow n|(r-r') $ or $ 0 \le r < n $ et $ -n<-r'\le 0\rightarrow r-r'=0 \rightarrow r=r'$
    \\On pourra aussi prouver l'inverse: $r=r' \rightarrow a-b=qn+r-(q'n-r')=(q-q')n+r-r' \rightarrow a-b=(q-q')n \rightarrow n|(a-b) \rightarrow a\equiv b(n)$
    \\\\Propriétés;
    \begin{enumerate}
        \item $a \equiv b(n)$ $b\equiv c(n) \rightarrow a\equiv c(n)$
        \item $a \equiv b(n)$ $c \equiv d(n) \rightarrow a+c\equiv b+d(n)$
        \item $a \equiv b(n)$ $c \equiv d(n) \rightarrow ac\equiv bd(n)$
        \item $a \equiv b(n)$ $k\in\mathbf{Z} \rightarrow ak\equiv bk(n)$
        \item $a \equiv b(n)$ $p\in\mathbf{N} \rightarrow a^p\equiv b^p(n)$
    \end{enumerate}


\subsection{Matrices et graphes}

\subsubsection{Notions de matrice}
    Définitions:
    \begin{enumerate}
        \item Une matrice est un tableau dimension lignes$\times$colonnes rempli de termes/coefficients.
        \item Une matrice ligne est une matrice dotée d'une unique ligne, et inversement pour une matrice colonne.
        \item Deux matrices sont égales si elles sont de même taille et de mêmes termes.
        \item Une matrice carrée d'ordre $p$ est une matrice de dimensions $p\times p$.
        \item Une matrice diagonale a tout les termes hors de sa diagonale (haut-gauche $\rightarrow$ bas-droite) de nuls.
        \item Une matrice est triangle supérieure ou inférieure si tout ses termes respectivement inférieurs ou supérieurs à sa diagonale sont nuls
        \item On note $I_p$ la matrice carrée d'ordre $p$ diagonale, dont les termes de la diagonale valent $1$.
        \item La transposée de la matrice $A$ de dimensions $n\times p$, notée $^TA$ (cette notation peut varier) de dimensions $p\times n$, obtenue en échangeant les colonnes et les lignes de $A$.
    \end{enumerate}

\subsubsection{Opérations sur les matrices}
\paragraph{Additions}
    \begin{enumerate}
        \item $A+B = a_{i,j}+b_{i,j}$
        \item $kA = ka_{i,j}$
        \item Cette addition est commutative.
        \item Les règles de distributivité et de factorisation s'appliquent.
    \end{enumerate}
    \hfill \break
\paragraph{Multiplication}
    \begin{enumerate}
        \item Pour effectuer le produit de $A$ par $B$, $A$ doit avoir autant de lignes que $B$ de colonnes.
        \item La matrice obtenue par le produit d'une matrice de taille $n \times p$ et une autre de $p \times q$ est de taille $n \times q$.
        \item $$A\times B=\sum_{k=1}^{p} a_{i,k} \times b_{k,j}$$
        \item Ce produit n'est pas commutatif.
        \item $A(n\times n) \times I_n = A$
    \end{enumerate}
\paragraph{Puissances}
    \begin{enumerate}
        \item $A^2 = AA$
        \item $A^0 = I$
    \end{enumerate}
\paragraph{Inverse d'une matrice($A^{-1}$)}:\\
    Une matrice est $A$ dite inversible ssi $\exists B, AB = BA = I$\\
    Dans le cas d'une matrice carré d'ordre 2 de la forme
        $\begin{pmatrix}
        a & b\\
        c & d
        \end{pmatrix}$
    , $A^-1 = \frac{1}{ad-bc}
        \begin{pmatrix}
        d & -b\\
        -c & a
        \end{pmatrix}$
    ssi $ad-bc \ne 0$.

\subsubsection{Systèmes grâce aux matrices}
    On peut écrire un système sous la forme de produit de la matrice $\mathbf{A}$ qui contient les coefficients, la matrice $\mathbf{X}$ qui contient les inconnues, et la matrice $\mathbf{B}$ qui contient les résultats. On trouve alors $\mathbf{AX} = \mathbf{B}$, donc $\mathbf{X} = \mathbf{A}^{-1}\mathbf{B}$.


\section{Boite à outils}


\subsection{Règles sur les inégalités}
    \begin{enumerate}
        \item Les additions et soustraction \emph{ne changent pas le sens}.
        \item Les multiplication et division \emph{ne changent le sens QUE SI leur produit ou dividende est négatif}.
        \item L'application d'une fonction \emph{ne change le sens QUE SI sa fonction est impaire}.
    \end{enumerate}


\subsection{Prouver une proposition}

\subsubsection{Par déduction}
    On part d'une hypothèse et par un raisonnement logique, on arrive à une conclusion.

\subsubsection{Par contre exemple}
    On trouve un cas qui contredit la proposition.

\subsubsection{Par l'absurde}
    On suppose P et arrive à une contradiction.

\subsubsection{Par contraposée}
    On prouve la contraposée pour prouvée la proposition

\subsubsection{Par disjonction de cas}
    Lorsque la proposition dépend d'un variable, on sépare le raisonnement suivant les valeurs de la dite variable.

\subsubsection{Par récurrence}
    On prouve qu'une proposition est vrai à partir d'un certain rang, en prouvant qu'elle est initialisée pour ce rang puis qu'elle est héréditaire.




\end{document}
